\chapter{Confronto sperimentale PIR vs mmWave}
\label{cap:confronto}

\begin{quote}\small
Questo è il capitolo centrale della tesi e risponde all'obiettivo O3. Presenta il
protocollo di misura, i risultati acquisiti e la loro discussione critica. Le
sezioni contrassegnate come \emph{in corso} indicano acquisizioni pianificate ma
non ancora completate alla data di questa stesura: la loro collocazione nel
capitolo è già definita, così come le metriche e i comandi con cui verranno
prodotte.
\end{quote}

\section{Protocollo sperimentale}
\label{sec:protocollo}

\subsection{Principi generali}

Il protocollo è stato definito \emph{prima} di iniziare le acquisizioni ed è
rimasto invariato, per una ragione precisa: la campagna completa richiede circa
dieci ore di acquisizioni, e una modifica del protocollo a metà strada le renderebbe
non confrontabili.

I principi adottati sono cinque.

\begin{enumerate}[leftmargin=1.8em,itemsep=4pt]
  \item \textbf{Cinque ripetizioni per scenario} (T01--T05). È ciò che permette di
        riportare media e deviazione standard, e quindi di rendere il confronto
        ``numerico'' e non anecdottico. Ogni valore riportato in questo capitolo
        nella forma $\mu \pm \sigma$ è calcolato su cinque trial indipendenti.
  \item \textbf{Acquisizione sincrona dei due sensori}. Radar e PIR sono letti nello
        stesso frame dal medesimo firmware, sulla stessa base dei tempi: il
        confronto non richiede allineamento a posteriori e non introduce errori di
        sincronizzazione.
  \item \textbf{Ground truth dichiarata}. Ogni acquisizione porta nei metadati la
        presenza reale e lo stato reale del soggetto, dichiarati dall'operatore da
        riga di comando. Per la ground truth statica per file questo è sufficiente;
        per le latenze si adotta la convenzione dell'evento a istante noto
        (\cref{sec:latenza}).
  \item \textbf{Setup fisso}. LD2410B e PIR affiancati sul bordo di un tavolo a
        circa 1~m di altezza, puntati verso l'area di prova, con la zona di prova
        interna al campo di \emph{entrambi} i sensori (il PIR copre un cono più
        ampio, $<$110\degree, del radar). Nessun oggetto in movimento nella stanza:
        nessun ventilatore, tende ferme, porta chiusa.
  \item \textbf{Alimentazione da alimentatore di rete} da almeno 1~A e non dalla
        porta USB del portatile. L'ESP32 con WiFi attivo e il LD2410B presentano
        picchi superiori a 400~mA, e una porta USB debole provoca brownout, cioè
        riavvii casuali che si presentano come bug del firmware.
\end{enumerate}

\subsection{Metriche}

Le metriche calcolate da \texttt{analisi/analizza\_test.py} sono le seguenti.

\begin{description}[leftmargin=1.4em,style=nextline,itemsep=3pt]
  \item[Falsi negativi (\%)] percentuale dei campioni con presenza reale in cui il
    sensore riporta assenza. È la metrica chiave dello scenario ``persona ferma''.
  \item[Falsi positivi (eventi/h)] numero di transizioni $0 \to 1$ per ora di
    osservazione a stanza vuota. Si conta per \emph{eventi} e non per campioni,
    perché un singolo falso positivo che dura dieci secondi non è dieci volte più
    grave di uno che dura un secondo.
  \item[Latenza di rilevamento] tempo fra l'evento reale e la prima transizione
    $0 \to 1$ del sensore.
  \item[Latenza di rilascio] tempo fra l'uscita del soggetto e la transizione
    stabile $1 \to 0$, con conteggio separato delle eventuali riaccensioni nella
    coda.
  \item[Accuratezza della distanza] scostamento fra distanza riportata dal radar e
    distanza reale misurata, con separazione fra dispersione \emph{fra} trial e
    dispersione \emph{entro} il trial.
\end{description}

\subsection{Registro delle sessioni}
\label{sec:registro}

Ogni sessione è annotata in un registro con data, ora, temperatura della stanza,
soggetto, alimentazione, test svolti e note. La temperatura è registrata perché il
PIR degrada al crescere della temperatura ambiente (\cref{sec:pir-temperatura}); le
sessioni di questa campagna si collocano fra 25 e 29\,\celsius, condizione
sfavorevole al PIR che va tenuta presente nella discussione (\cref{sec:limiti}).

Le acquisizioni completate sono riassunte nella \cref{tab:sessioni}.

\begin{table}[htbp]
\centering
\caption{Sessioni di acquisizione completate.}
\label{tab:sessioni}
\small
\begin{tabularx}{\textwidth}{@{}llrX@{}}
\toprule
\textbf{Data} & \textbf{Scenario} & \textbf{Durata} & \textbf{Note} \\
\midrule
18/08/2026 & validazione logger & 134~s & 673 campioni, 5,00~Hz esatti \\
18/08/2026 & pilota PIR vs radar & 339,6~s & 1699 campioni, PIR su GPIO34 \\
18/08/2026 & pilota respiro & 91~s & soggetto a 40~cm, nessuna ground truth \\
18/08/2026 & stanza vuota (giorno) & 1742~s & 8711 campioni, 29\,\celsius \\
18/08/2026 & fermo seduto T01--T05 & 5 $\times$ 302~s & 7554 campioni, 2,30~m, 27\,\celsius \\
20/08/2026 & distanza 1--5~m, 25 trial & 25 $\times$ 62~s & cammino sul posto \\
20--21/08/2026 & stanza vuota (notte) & 23640,6~s & 118\,204 campioni, 6,57~h \\
\bottomrule
\end{tabularx}
\end{table}

\paragraph{Nota su una sessione troncata.} La sessione notturna era programmata per
sette ore ed è terminata a 6,57~h perché un aggiornamento automatico di Windows ha
riavviato il computer. I dati restano validi: l'intervallo di campionamento è di
200~ms esatti su tutti i 118\,203 intervalli, senza buchi né campioni persi ---
l'acquisizione è semplicemente troncata. L'episodio è però istruttivo e ha prodotto
una regola di protocollo: per le acquisizioni oltre l'ora, sospendere gli
aggiornamenti di Windows e verificare che sospensione e ibernazione siano disattivate.

\section{Falsi positivi: stanza vuota}
\label{sec:falsi-positivi}

\paragraph{Procedura.} Acquisizione avviata con l'operatore che esce immediatamente
dalla stanza; scarto dei primi 60~s (stabilizzazione del PIR e uscita
dell'operatore); conteggio delle transizioni $0 \to 1$ sul resto.

\paragraph{Risultato.} Su un totale di \textbf{7,05~h} di stanza vuota
(6,57~h notturni più 0,48~h diurni), \textbf{né il radar né il PIR hanno prodotto
alcun falso positivo}. Nella sessione notturna gli unici 96 campioni con presenza
sono compresi fra $t=0$ e $t=19$~s e corrispondono all'operatore che esce dalla
stanza.

Con zero eventi osservati non è corretto scrivere ``0 eventi/h'': si riporta il
limite superiore, che dipende dal tempo di osservazione. Applicando la regola del
tre ($3/T$, che approssima il limite superiore al 95\,\% di confidenza per un
processo di Poisson con zero eventi osservati):

\begin{equation}
\lambda_{95\%} \leq \frac{3}{7{,}05~\text{h}} = 0{,}43~\text{eventi/h}
\label{eq:regola-tre}
\end{equation}

Il risultato da riportare è quindi: \textbf{falsi positivi $\leq 0{,}43$ eventi/h
per entrambi i sensori, al 95\,\% di confidenza}. Il valore è coerente con l'ordine
di grandezza indicato dalla letteratura tecnica per i moduli mmWave consumer
($<$0,05 eventi/h), che però richiederebbe una sessione di decine di ore per essere
confermato con questa metodologia. Questo è il modo corretto di riportare
un'assenza di eventi, e la principale limitazione del risultato è il tempo di
osservazione, non il sensore.

\paragraph{Nota sulla temperatura.} Le sessioni sono state svolte a 29\,\celsius. È
una condizione \emph{sfavorevole} al PIR per il rilevamento (minore contrasto
termico) ma \emph{favorevole} sui falsi positivi: un ambiente caldo e uniforme
produce meno gradienti spuri. Il risultato ``zero falsi positivi PIR'' va quindi
letto come non pessimistico.

\section{Il risultato centrale: la persona ferma}
\label{sec:persona-ferma}

\paragraph{Procedura.} Soggetto seduto immobile a 2,30~m dalla coppia di sensori,
respirazione spontanea, nessun movimento volontario. Cinque trial da 302~s utili
ciascuno dopo lo scarto dei primi 30~s. Temperatura della stanza 27\,\celsius.
Totale: 7554 campioni con presenza reale.

\paragraph{Risultato.} I valori misurati sono riportati nella
\cref{tab:fermo-seduto}.

\begin{table}[htbp]
\centering
\caption{Scenario ``persona ferma seduta a 2,30~m''. Falsi negativi sui 7554
campioni con presenza reale, su cinque trial indipendenti da 302~s.}
\label{tab:fermo-seduto}
\small
\begin{tabular}{@{}lrr@{}}
\toprule
\textbf{Metrica} & \textbf{PIR \pir} & \textbf{Radar \ldb} \\
\midrule
Falsi negativi (media $\pm$ dev.\ std.\ su 5 trial) & $99{,}78 \pm 0{,}49\,\%$ &
  $\mathbf{0{,}00 \pm 0{,}00\,\%}$ \\
Trial con falsi negativi al 100\,\% & 4 su 5 & 0 su 5 \\
Campioni con presenza rilevata & $\sim$17 su 7554 & 7554 su 7554 \\
\bottomrule
\end{tabular}
\end{table}

Il dato va letto con attenzione al suo significato operativo: in quattro trial su
cinque il PIR ha prodotto \textbf{zero rilevamenti in cinque minuti} con una
persona viva seduta a 2,3~m di distanza, in campo visivo libero. Il radar ha
rilevato la presenza nel 100\,\% dei campioni in tutti e cinque i trial.

Questo è il risultato che risponde alla domanda di ricerca principale. Non è un
esito sorprendente --- è precisamente ciò che la fisica del piroelettrico prevede
(\cref{sec:pir-principio}) --- ma la tesi lo \emph{quantifica}, con deviazione
standard su cinque ripetizioni, nella condizione che interessa il progetto UPRISE.

Un risultato analogo, ottenuto in una sessione precedente con procedura meno
controllata, conferma il quadro: in un tratto di \textbf{93,2~s di immobilità
continuativa} a circa 2,9~m, il radar ha rilevato la presenza nel 100\,\% dei
campioni e il PIR nello 0,2\,\% (un solo campione su 466).

\subsection{Come va interpretato il valore percentuale}
\label{sec:interpretazione-pir}

Il numero ``99,78\,\% di falsi negativi'' richiede una precisazione onesta, perché
dipende in parte da un parametro sotto il nostro controllo: il trimmer del tempo di
ritenuta, che abbiamo tenuto al minimo (\cref{sec:hcsr501}).

Nella sessione pilota il PIR ha prodotto \textbf{14 impulsi di durata costante
compresa fra 3,4 e 3,6~s} invece di un segnale continuo. Su due sessioni sono stati
misurati in totale \textbf{30 impulsi, tutti compresi fra 3,4 e 3,8~s}. Questa è la
prova quantitativa di un fatto qualitativo: \textbf{l'uscita del PIR è un
monostabile a durata fissa}, indipendente da cosa fa la persona. Il PIR non misura
la presenza né la durata del movimento: segnala eventi.

La conseguenza è che, con la ritenuta regolata al massimo, il PIR
\emph{sembrerebbe} rilevare molto di più --- ma starebbe soltanto trattenendo
l'ultimo evento per minuti. La percentuale di campioni rilevati non è quindi una
proprietà robusta del sensore.

Il dato robusto e non contestabile è l'altro, e va riportato come tale: nei 93~s di
immobilità continuativa e nei cinque trial da cinque minuti, \textbf{il PIR non ha
prodotto alcun evento di movimento} (con le eccezioni contate sopra), perché non
c'era alcun transito di flusso infrarosso da rilevare. La percentuale di campioni è
il modo di quantificarlo, il numero di eventi è il modo di dimostrarlo.

\section{Il PIR sbaglia anche con la persona in movimento}
\label{sec:pir-movimento}

Questo è il risultato meno atteso della campagna, e con il maggiore impatto sul
contesto applicativo.

\paragraph{Procedura.} Soggetto che cammina \emph{sul posto} a distanza fissa dalla
coppia di sensori --- movimento continuo, ma senza spostamento nello spazio. Cinque
distanze (1, 2, 3, 4, 5~m), cinque trial per distanza, 62~s utili ciascuno: 25
trial in totale.

\paragraph{Risultato.} A 1~m, condizione teoricamente ideale per un PIR (bersaglio
vicino, movimento continuo, campo libero), il PIR ha registrato
$\mathbf{80{,}46 \pm 6{,}22\,\%}$ di falsi negativi. In 62~s di cammino continuo ha
emesso da \textbf{1 a 4 impulsi} per trial, di durata 3,6--3,8~s. A 2, 3, 4 e 5~m il
tasso di falsi negativi è \textbf{100,0\,\%} su tutti e venti i trial: nessun
rilevamento. Il radar ha rilevato la presenza nel 100\,\% dei campioni a tutte e
cinque le distanze.

\paragraph{Spiegazione.} Il risultato non è un'anomalia e si spiega esattamente con
il meccanismo della lente di Fresnel (\cref{sec:fresnel}): il piroelettrico
risponde al \emph{transito} del flusso infrarosso fra le zone della lente, non al
movimento in sé. Chi si muove \emph{sul posto} non attraversa le zone, quindi non
genera transizioni --- e per il sensore è indistinguibile da un corpo fermo.

\paragraph{Rilevanza per UPRISE.} Questo risultato è più importante di quanto la
sua sobrietà suggerisca. Una persona intrappolata sotto un arredo \emph{si muove sul
posto}: si agita, cerca di liberarsi, sposta le braccia, chiama aiuto. Non attraversa
la stanza. È il caso peggiore per un PIR, e la campagna mostra che il PIR fallisce
in questo caso quasi allo stesso modo in cui fallisce sulla persona completamente
immobile ($80\,\%$ contro $\sim\!100\,\%$ di falsi negativi, e $100\,\%$ oltre i 2~m).

Detto in altri termini: l'ipotesi implicita ``il PIR non vede chi è fermo ma vede
chi si muove'' è, nello scenario di interesse, falsa.

\section{Accuratezza della misura di distanza}
\label{sec:distanza}

Gli stessi 25 trial forniscono la caratterizzazione della misura di distanza del
radar, riportata nella \cref{tab:distanza}.

\begin{table}[htbp]
\centering
\caption{Accuratezza della distanza riportata dal \ldb, su 25 trial (5 distanze
$\times$ 5 ripetizioni). L'energia è il valore medio del bersaglio, normalizzato
0--100.}
\label{tab:distanza}
\small
\begin{tabular}{@{}rrrrr@{}}
\toprule
\textbf{Distanza reale} & \textbf{Errore grezzo} & \textbf{Errore relativo} &
\textbf{Dispersione entro trial} & \textbf{Energia media} \\
\midrule
1~m & $-0{,}5$~cm & $-0{,}5\,\%$ & 10--26~cm & 99 \\
2~m & --- & --- & 10--26~cm & 85 \\
3~m & --- & --- & 10--26~cm & 54 \\
4~m & --- & --- & 10--26~cm & 35 \\
5~m & $+18{,}2$~cm & $+3{,}6\,\%$ & 10--26~cm & 28 \\
\bottomrule
\end{tabular}
\end{table}

\todo{completare le colonne intermedie della \cref{tab:distanza} con i valori
puntuali per 2, 3 e 4~m e la dispersione entro trial per singola distanza,
rileggendoli dall'output di \texttt{analizza\_test.py}}

\subsection{Il modello di errore}

La regressione lineare fra distanza misurata e distanza reale, sui 25 trial, dà:

\begin{equation}
d_{\text{misurata}} = 1{,}0381 \cdot d_{\text{reale}} - 1{,}32~\text{cm},
\qquad R^2 = 0{,}99965
\label{eq:regressione}
\end{equation}

Il risultato si legge in due parti, ed è la lettura --- non il numero grezzo --- il
contributo di questa sezione.

\begin{itemize}[leftmargin=1.4em,itemsep=3pt]
  \item Il radar è \textbf{estremamente lineare}: $R^2 = 0{,}99965$ e residui dalla
        retta non superiori a 4,9~cm su tutto il range. Non c'è alcuna non linearità
        apprezzabile fra 1 e 5~m.
  \item L'errore è quasi interamente un \textbf{errore di scala del $+3{,}81\,\%$},
        con offset praticamente nullo ($-1{,}32$~cm). Questo spiega perché l'errore
        grezzo cresce con la distanza (da $-0{,}5$~cm a 1~m fino a $+18{,}2$~cm a
        5~m) ed è, soprattutto, \textbf{correggibile con un solo coefficiente}:
        dividendo la lettura per 1,0381 l'errore residuo scende sotto i 5~cm a tutte
        le distanze.
\end{itemize}

\paragraph{Limite dichiarato sulla causa.} I dati raccolti \emph{non} permettono di
attribuire l'origine del $+3{,}81\,\%$. Può trattarsi della calibrazione interna del
modulo oppure del riferimento con cui sono state segnate le distanze sul pavimento.
Per discriminare servirebbe un riferimento di distanza indipendente, per esempio un
distanziometro laser. La distinzione è rilevante: nel primo caso il coefficiente
correttivo è una proprietà del modulo e vale in generale, nel secondo è un artefatto
del nostro setup.

\subsection{Dispersione entro il trial}

La dispersione della distanza \emph{entro} un singolo trial è risultata di
10--26~cm nello scenario di cammino sul posto, contro
$\mathbf{1{,}48 \pm 0{,}77}$~\textbf{cm} nello scenario da seduto immobile
(\cref{sec:persona-ferma}). Il confronto è informativo: la dispersione maggiore non
è rumore del sensore, è l'oscillazione reale del busto di chi cammina sul posto. Il
radar la sta misurando correttamente.

Nello scenario da seduto emerge una scomposizione utile: la variabilità \emph{fra}
trial della distanza media è di $\pm 2{,}38$~cm (media $229{,}78$~cm su una distanza
reale di 230~cm), mentre la stabilità \emph{entro} il trial è di $1{,}48$~cm. La
variabilità di come il soggetto si siede ($\sim$2,4~cm) è quindi maggiore di quella
del sensore ($\sim$1,5~cm): il fattore limitante dell'accuratezza in questo scenario
è il soggetto, non lo strumento.

\subsection{Decadimento dell'energia con la distanza}

L'energia media del bersaglio decade da 99 a 1~m fino a 28 a 5~m
(\cref{tab:distanza}). Il decadimento è \emph{molto} più lento di quanto
prevederebbe l'equazione del radar, che per un bersaglio puntiforme dà una
dipendenza $1/D^4$ della potenza ricevuta.

Va quindi evitata una tentazione interpretativa: \textbf{questo valore non è una
potenza e non va letto con l'equazione del radar}. È un indicatore normalizzato
0--100 prodotto da un'elaborazione interna non documentata, con ogni probabilità
compressiva. Se ne può usare la monotonia (più lontano $\to$ meno energia) ma non il
valore assoluto in termini fisici.

L'osservazione ha però un uso pratico: l'extrapolazione della tendenza porta a un
valore di circa 27 a 6~m, contro una soglia di rilevamento di 15 sul gate
corrispondente (\cref{tab:soglie}). Questo indica che il sensore avrebbe funzionato
anche a 6~m.

\paragraph{Perimetro sperimentale.} Il range coperto è \textbf{1--5~m} e non 6~m: è
un limite della stanza disponibile, non del sensore, e va dichiarato come tale. La
retta di regressione è costruita su cinque punti e 25 trial. Per lo scenario UPRISE
la limitazione è comunque irrilevante, perché la distanza di interesse è
\emph{inferiore al metro} --- la persona sotto il banco --- che è coperta dallo
scenario di \cref{sec:sotto-banco}.

\section{Latenze di rilevamento e di rilascio}
\label{sec:latenza}

\emph{Acquisizione in corso.}

\paragraph{Metodo previsto.} Poiché la ground truth è statica per file, la latenza
si misura con la convenzione dell'\textbf{evento a istante noto}: l'acquisizione
parte con la stanza vuota e il soggetto entra esattamente a $t = 10$~s, con
riferimento condiviso fra operatore e acquisizione. Lo script di analisi riceve
l'istante dell'evento come parametro (\texttt{-{}-event-time} per il rilevamento,
\texttt{-{}-release-time} per il rilascio) e calcola la latenza come distanza
temporale fra l'evento e la prima transizione stabile.

\paragraph{Attese e vincoli noti.} Tre elementi sono già stabiliti e vanno riportati
con i risultati.

\begin{itemize}[leftmargin=1.4em,itemsep=2pt]
  \item La latenza di rilascio del radar \textbf{non è un difetto}: è il
        \emph{no-one duration}, impostato a 5~s di fabbrica, che nella sessione di
        validazione ha prodotto una coda osservata di circa 10~s con un bersaglio
        fantasma in decadimento a 5,2~m (\cref{sec:coda}). Va misurata come tale e
        non conteggiata fra i falsi positivi.
  \item Il PIR ha un \textbf{block time di 2,5~s} dopo ogni rilascio, durante il
        quale è cieco per costruzione. Una latenza di rilevamento misurata subito
        dopo un rilascio è quindi contaminata da questo effetto e va segnalata.
  \item La direzione di ingresso conta, e in senso opposto per i due sensori: il PIR
        favorisce l'ingresso \emph{trasversale}, il radar quello \emph{radiale}
        (\cref{sec:fresnel}). La direzione va annotata per ogni trial e i due casi
        vanno riportati separatamente, altrimenti il confronto è ambiguo.
\end{itemize}

\section{Lo scenario UPRISE: persona sotto il banco}
\label{sec:sotto-banco}

\emph{Acquisizione in corso.}

È lo scenario che replica la condizione d'interesse del progetto: soggetto
accovacciato sotto un piano, a distanza inferiore al metro dal sensore. Rispetto
agli scenari precedenti presenta tre elementi nuovi, tutti già identificati:

\begin{enumerate}[leftmargin=1.8em,itemsep=3pt]
  \item \textbf{distanza ravvicinata}, quindi rischio di saturazione dei canali di
        energia (\cref{sec:saturazione});
  \item \textbf{canali stazionari per-gate non disponibili} sotto 1,5~m
        (\cref{sec:senergy-zero}): l'analisi dovrà usare i canali in movimento,
        e questa verifica è esplicitamente parte dello scenario;
  \item \textbf{presenza del piano dell'arredo} fra sensore e soggetto, a seconda
        del punto di montaggio. Nel banco UPRISE il piano è rinforzato con
        \emph{lamiera forata}: il metallo è il materiale peggiore per un radar
        (\cref{sec:ostacoli}), e la posizione di montaggio del sensore rispetto alla
        lamiera è una domanda aperta di progetto, da porre al relatore.
\end{enumerate}

\section{Penetrazione degli ostacoli}
\label{sec:ostacoli}

\emph{Acquisizione in corso.}

\paragraph{Metodo previsto.} Ripetizione dello scenario ``persona ferma'' con
l'interposizione di un pannello fra sensore e soggetto, per i materiali disponibili:
cartongesso, legno, vetro, plastica. Metrica: tasso di falsi negativi e riduzione
dell'energia del bersaglio rispetto al caso senza ostacolo, a parità di distanza.

\paragraph{Attese da letteratura.} Il radar a 24~GHz attraversa cartongesso, legno,
vetro e plastica con attenuazione contenuta, mentre è bloccato da metallo e da
cemento armato di spessore rilevante. Il PIR è invece bloccato da \emph{tutti} i
materiali elencati, vetro compreso: la finestra ottica 5--14~\textmu m del corpo
umano non attraversa il vetro comune. Quest'ultimo punto è un risultato atteso ma
importante per il confronto, perché rende il PIR inutilizzabile per qualunque
montaggio incassato o protetto.

\section{Il rilevamento del respiro}
\label{sec:respiro}

\subsection{Metodo}

Il metodo segue quanto descritto in \cref{sec:respiro-teoria}: si campiona a 5~Hz
un canale di energia per-gate in engineering mode, si applica la FFT alla serie
temporale e si cerca il picco in banda 0,1--0,5~Hz. La frequenza del picco,
moltiplicata per 60, dà la stima degli atti al minuto.

Il contributo metodologico di questa tesi è nella \textbf{selezione del canale}. Il
LD2410B mette a disposizione venti serie temporali di energia (due di bersaglio più
diciotto per-gate) e non tutte sono utilizzabili: alcune sono sature
(\cref{sec:saturazione}), altre identicamente nulle (\cref{sec:senergy-zero}), altre
piatte perché lontane dal bersaglio. La modalità \texttt{-{}-scan} dello script di
analisi esamina tutti i canali, scarta quelli saturi e quelli piatti secondo criteri
dichiarati, ordina i rimanenti per rapporto segnale-rumore e riporta la
\textbf{mediana delle stime concordi}. La concordanza fra canali indipendenti è
l'unico controllo di plausibilità disponibile in assenza di ground truth.

\subsection{Prima prova: soggetto a 40~cm}

\paragraph{Risultato.} Soggetto fermo a circa 40~cm, 91~s di acquisizione,
respirazione spontanea. \textbf{Sette canali di energia in movimento, indipendenti
fra loro, concordano su 0,264~Hz, cioè 15,8 atti al minuto}, con rapporto
segnale-rumore fino a 12,8 sul canale \texttt{menergy\_gate2}. Il valore è
fisiologicamente plausibile per un adulto a riposo.

\paragraph{Perché non è una misura valida.} Va detto con chiarezza, perché è una
distinzione che qualifica il lavoro: questa prova dimostra la \emph{fattibilità} ma
non è una misura utilizzabile come risultato. Le ragioni sono tre:
\emph{(i)} manca la ground truth, cioè non si conosce il ritmo respiratorio reale e
quindi non si può calcolare un errore; \emph{(ii)} l'89\,\% dei campioni sui canali
in movimento è saturo; \emph{(iii)} 40~cm non è una distanza rappresentativa di
alcuno scenario d'uso.

\paragraph{I canali stazionari sono inutilizzabili a questa distanza.} I gate 2 e 3
stazionari sono saturi al 100\,\%, i gate 0 e 1 sono identicamente nulli
(\cref{sec:senergy-zero}) e i gate 4--8 forniscono stime fra 6,6 e 9,2 atti al
minuto, in disaccordo fra loro e con il canale in movimento. Il valore 0,132~Hz che
appare su alcuni di questi canali è esattamente la metà di 0,264~Hz: si tratta con
ogni probabilità di una subarmonica, cioè di un artefatto dell'analisi, non di un
segnale.

\subsection{Seconda prova: soggetto a 2,31~m}

La stessa analisi è stata applicata ai cinque trial dello scenario ``persona ferma''
(\cref{sec:persona-ferma}), a 2,31~m e per 302~s ciascuno. Il quadro si ripete, ma
in condizioni migliori.

\begin{itemize}[leftmargin=1.4em,itemsep=3pt]
  \item I canali \textbf{in movimento} dei gate 2--8 concordano su circa 0,30~Hz,
        cioè \textbf{18,2 atti al minuto}, valore plausibile per un adulto seduto.
        Le stime per i cinque trial sono 18,2 / 18,9 / 18,8 / 20,0 / 20,2 atti al
        minuto, con 7--8 canali concordi per trial.
  \item Questi canali \textbf{non sono saturi} a questa distanza: la percentuale di
        campioni a fondo scala è fra 0 e 1\,\%. È una buona notizia di natura
        pratica, perché significa che la prova definitiva del respiro può essere
        svolta a distanza realistica senza accorgimenti geometrici particolari.
  \item I canali \textbf{stazionari} danno 6,5--7,8 atti al minuto, troppo pochi per
        un adulto a riposo, e i gate prossimi al bersaglio sono saturi (gate~3 al
        100\,\%, gate~4 al 97\,\%).
\end{itemize}

\paragraph{Conclusione operativa.} Due sessioni a distanze molto diverse (40~cm e
2,31~m) dicono la stessa cosa: \textbf{per il respiro vanno usati i canali in
movimento}. Il valore intorno a 7 atti al minuto dei canali stazionari è
verosimilmente un artefatto del filtraggio interno del canale e non respirazione ---
ma, in assenza di ground truth, questa resta un'ipotesi e va presentata come tale.
A deciderlo è la prova con metronomo descritta di seguito.

\subsection{La prova definitiva: respiro a ritmo imposto}
\label{sec:respiro-metronomo}

\emph{Acquisizione in corso.}

La misura valida per la tesi richiede una ground truth, e il modo più semplice di
ottenerla è imporre il ritmo: il soggetto respira seguendo un metronomo a
frequenza nota (per esempio 10, 15 e 20 atti al minuto), a distanza $\geq 1{,}5$~m,
per cinque trial per ritmo. Il confronto fra stima e ritmo imposto è la vera prova,
e permette di riportare un errore assoluto in atti al minuto invece di una
concordanza fra canali.

È inoltre necessario un \textbf{controllo negativo}: la stessa analisi applicata a
una registrazione di stanza vuota deve \emph{non} produrre un picco in banda
respiratoria. Senza questo controllo, un picco nello spettro non dimostra nulla, e
la banda 0,1--0,5~Hz è abbastanza larga da ospitare artefatti.

\section{Confronto fra i due moduli radar}
\label{sec:ld2420-risultati}

\emph{Attività parziale.}

Il LD2420 è stato collegato e verificato (\cref{sec:ld2420}): risponde ai movimenti,
riporta presenza e un valore di ``range'' in modalità ASCII a 115200 baud. In questa
modalità però non fornisce energia per-gate né distingue movimento e stazionario, e
l'unità del campo ``range'' non è documentata. Il confronto quantitativo con il
LD2410B richiede quindi la riconfigurazione del modulo in modalità binaria.

Alla data di questa stesura il confronto è limitato a un raffronto di specifiche
(\cref{tab:confronto-moduli}) e alla constatazione qualitativa che il LD2420
rileva presenza a distanze superiori. Il posizionamento dei due moduli resta quello
indicato dalle specifiche: il LD2410B è il modulo adatto al caso UPRISE (dettaglio
per-gate, respiro, distanze brevi), il LD2420 a una copertura di ambiente più ampia.

\section{Discussione}
\label{sec:discussione}

\subsection{Dove il mmWave vince}

Il quadro che emerge dai dati acquisiti è netto e si può riassumere in quattro
punti.

\begin{enumerate}[leftmargin=1.8em,itemsep=3pt]
  \item \textbf{Sulla persona ferma il divario non è graduale, è categorico}:
        $0{,}00\,\%$ contro $99{,}78\,\%$ di falsi negativi. Non si tratta di un
        sensore migliore di un altro, ma di un sensore che misura una grandezza che
        l'altro non ha accesso a misurare.
  \item \textbf{Il divario si estende alla persona in movimento sul posto}, che è il
        caso applicativo reale: $100\,\%$ di falsi negativi PIR oltre i 2~m contro
        $0\,\%$ del radar.
  \item \textbf{Il radar aggiunge informazione, non solo affidabilità}: la distanza
        è misurata con linearità $R^2 = 0{,}99965$ e un errore correggibile con un
        coefficiente; l'energia per-gate consente l'analisi del respiro e l'indice
        di vitalità (\cref{cap:vitalita}), cose che con un bit non sono
        rappresentabili.
  \item \textbf{Sui falsi positivi i due sensori sono equivalenti}: entrambi
        $\leq 0{,}43$ eventi/h su 7~h di osservazione. Il vantaggio del radar non si
        paga in affidabilità.
\end{enumerate}

\subsection{Dove il mmWave non vince}

Per equilibrio vanno riportati anche i punti a sfavore, che sono reali.

\begin{itemize}[leftmargin=1.4em,itemsep=3pt]
  \item \textbf{Consumo}: tre ordini di grandezza in più (80~mA contro 50~\textmu A).
        È il vincolo che, da solo, determina l'architettura raccomandata nel
        \cref{cap:conclusioni}, e il \cref{cap:consumi} lo quantifica.
  \item \textbf{Coda di presenza}: il radar mantiene la presenza per circa 10~s dopo
        l'uscita. Per un'applicazione di illuminazione automatica sarebbe un
        vantaggio, per una misura di occupazione istantanea è un ritardo.
  \item \textbf{Saturazione a distanza ravvicinata}: proprio nello scenario di
        interesse (sotto il banco) i canali più informativi rischiano di saturare,
        e i canali stazionari per-gate non sono disponibili sotto 1,5~m. Sono
        vincoli di progetto, non insormontabili, ma vanno gestiti.
  \item \textbf{Bersaglio singolo}: il LD2410B non dispone di un array di antenne e
        non distingue due persone alla stessa distanza. Non fornisce informazione
        angolare.
  \item \textbf{Complessità}: richiede UART a 256000 baud, gestione di un protocollo
        a frame, e --- per usarne i dati veri --- engineering mode e analisi
        numerica. Il PIR richiede un piedino digitale.
\end{itemize}

\subsection{Limiti del setup sperimentale}
\label{sec:limiti}

I limiti da dichiarare esplicitamente sono i seguenti.

\begin{description}[leftmargin=1.4em,style=nextline,itemsep=3pt]
  \item[Numero di soggetti] Le acquisizioni sono state svolte con un solo soggetto
    (indicato come S1 nel registro). La corporatura influisce sull'energia riflessa
    e, nel caso del PIR, sull'emissione infrarossa: la ripetizione degli scenari
    chiave con un secondo soggetto è la prima estensione da fare.
  \item[Ambiente] Stanza domestica, non un'aula scolastica: dimensioni, materiali,
    riflessioni multiple e presenza di arredi metallici sono diversi. Il
    trasferimento dei risultati al dimostratore reale va verificato.
  \item[Temperatura] Fra 25 e 29\,\celsius. Condizione sfavorevole al PIR in
    rilevamento: un confronto a 18--20\,\celsius\ è un'estensione a costo quasi
    nullo che rafforzerebbe il risultato, mostrando quanto della cecità sul fermo
    sia strutturale e quanto termico. La previsione, dalla fisica del principio, è
    che la cecità sul fermo resti totale a qualunque temperatura.
  \item[Range] 1--5~m per la distanza, per limite della stanza. Il caso UPRISE è
    sotto il metro e non è interessato.
  \item[Un solo esemplare per modello] Nessuna stima della variabilità fra esemplari.
    Il coefficiente di scala del $+3{,}81\,\%$ potrebbe essere specifico
    dell'esemplare in prova.
  \item[Assenza di strumentazione di misura] I consumi sono da datasheet
    (\cref{cap:consumi}) e le distanze sono riferite a marcature su pavimento, non a
    un distanziometro certificato.
\end{description}

Nessuno di questi limiti mette in discussione il risultato centrale, che è
robusto per la natura stessa del divario misurato: la cecità del PIR sulla persona
ferma è una conseguenza del principio fisico, e i dati la confermano nella
condizione di interesse con cinque ripetizioni indipendenti.
